% =============================================================================
% KARMA paper -- supplementary appendix (extended)
% =============================================================================
% These sections were trimmed from the LNCS submission (input/appendix.tex in
% the thirag-paper repo) to fit the page budget. They are preserved here in
% the dbpedia_qa codebase as supplementary material for reviewers and readers
% who want the full evidence: per-benchmark routing counts, two narrated
% execution paths, and a complete end-to-end LLM-prompt trace of one query
% through every operator of Algorithm~1 (\textsc{ThiRAG}).
%
% Cross-references in this file (Algorithm~\ref{alg:thirag},
% Table~\ref{tab:trace}, Figure~\ref{fig:routing-cross-benchmark}, citation
% keys) point into the main paper PDF.
% =============================================================================

\section{Cross-benchmark routing breakdown}
\label{sec:routing-breakdown}

Figure~\ref{fig:routing-cross-benchmark} in the main text shows
the Q-level and Sub-Q-level routing distributions;
Table~\ref{tab:routing-cross-benchmark-appendix} gives the
underlying routing counts.

\begin{table}[t]
\caption{Per-benchmark routing statistics over 1{,}000 questions each.
Direct SPARQL: whole question answered by Text2SPARQL without
decomposition. Sub-Qs: total sub-questions produced when decomposition
fires. Sub-SPARQL\%: percentage of generated sub-questions resolved
symbolically. Any-SPARQL\%: fraction of original questions in which at
least one sub-question hit the SPARQL path.
Routing distributions on 2WikiMHQA and MuSiQue are stable
across runs (any-SPARQL drift $\le 2.6\%$); HotpotQA is
more variable ($15.4\%$), reflecting LLM stochasticity in
the coverage check when KG support is sparse.}
\label{tab:routing-cross-benchmark-appendix}
\centering
\footnotesize
\setlength{\tabcolsep}{4pt}
\begin{tabular}{lrrr}
\toprule
\textbf{Metric} & \textbf{2WikiMHQA} & \textbf{HotpotQA} & \textbf{MuSiQue} \\
\midrule
Direct SPARQL (\%)            & 29.5  & 6.1   & 2.4  \\
Decomposition fired (\%)      & 70.1  & 88.7  & 96.5 \\
\quad Sub-questions generated & 1544  & 1807  & 2508 \\
\quad Sub-SPARQL (\%)         & 40.5  & 7.6   & 18.0 \\
\midrule
\textbf{Any-SPARQL touch (\%)} & \textbf{70.0} & \textbf{17.4} & \textbf{42.6} \\
\bottomrule
\end{tabular}
\end{table}

\subsection{Running examples: two execution paths}
\label{app:running-examples}

We illustrate the two execution paths of Algorithm~\ref{alg:thirag}
on questions from 2WikiMultiHopQA over DBpedia.  Path~A exercises
the KG-only fast path (lines~3--6 of Algorithm~\ref{alg:thirag});
Path~B exercises the decompose-and-mixed-routing branch
(lines~7--19).  The step-by-step trace in the main text
(Table~\ref{tab:trace}) takes a third question end-to-end through
every operator.

\smallskip\noindent\emph{Path~A --- KG-only fast path (lines~3--6).}\enspace
Consider $Q{=}$``Where was the director of film Thomas Jefferson (Film) born?''
$\textsc{EntityLink}$ maps \emph{Thomas Jefferson} to
\texttt{dbr:Thomas\_Jefferson\_(film)} (class \texttt{dbo:Film}).
$\textsc{SchemaProject}$ retrieves relevant properties
\texttt{dbo:director}, \texttt{dbo:birthPlace}; the projected schema
covers the full 2-hop chain required by~$Q$.
The coverage check at line~3 succeeds, so $\textsc{Text2SPARQL}$
generates the query
\texttt{\small SELECT ?bp WHERE \{ dbr:Thomas\_Jefferson\_(film)
dbo:director/dbo:birthPlace ?bp \}}, which returns \emph{Brooklyn}.
The answer is returned at line~5 without decomposition.

\smallskip\noindent\emph{Path~B --- Decompose + mixed routing
  (lines~7--19).}\enspace
Consider $Q{=}$``What nationality is Matilda of Frisia's husband?''
$\textsc{EntityLink}$ maps \emph{Matilda of Frisia} to
\texttt{dbr:Matilda\_of\_Frisia} (class \texttt{dbo:Person}).
$\textsc{SchemaProject}$ finds that the \texttt{dbo:spouse} property
\emph{is} present in this entity's local schema, but the required
second-hop property (\texttt{nationality} of the spouse) is not
directly accessible.  The coverage check at line~3 fails: the
projected schema cannot answer~$Q$ in a single query.
The system decomposes~$Q$ into two subproblems:
$p_1{=}$``Who is Matilda of Frisia's husband?'' and
$p_2{=}$``What nationality is [$p_1$]?'' with dependency $p_2
\!\to\! p_1$.
For~$p_1$, schema re-scoping confirms that the \texttt{spouse}
property is present ($\mathcal{S}_1$ covers~$\hat{p}_1$), so
$\textsc{Text2SPARQL}$ returns \emph{Henry~I of France}.
For~$p_2$, the resolved entity \texttt{dbr:Henry\_I\_of\_France}
does not expose a usable \texttt{nationality} property in the
projected schema, so $\textsc{TextRetrieve}$ retrieves passages
from Wikipedia confirming that Henry~I was \emph{French}.
$\textsc{Synthesize}$ produces the final answer: \emph{French}.

\smallskip\noindent
Path~A shows how the direct SPARQL check avoids unnecessary
decomposition when the KG is complete.  Path~B shows the text
fallback: the system decomposes the question, resolves the first
subproblem symbolically via SPARQL, and falls back to text for the
second.  The per-subproblem routing is decisive.
On the same question, LinearRAG~\cite{Zhuang2025} retrieves passages
that describe Henry~I as ``\emph{king of the Franks}'' and mention
``\emph{Frankish} domination'' of Frisia in the same context window;
it produces the answer \emph{Frankish} (marked incorrect by both
judges).  The confusion arises because LinearRAG attempts to answer
the two-hop question in a single retrieval pass: the passages that
mention Matilda's husband also use the medieval title ``king of the
Franks,'' and the model conflates this historical designation with
nationality.  \sysname avoids this failure by decomposing the
question and grounding the intermediate entity through SPARQL:
$p_1$ returns the unambiguous resource
\texttt{dbr:Henry\_I\_of\_France}, and $p_2$ retrieves text
specifically about Henry~I, where the nationality \emph{French}
is stated directly.

\subsection{End-to-End Trace with Full LLM Prompts for Mixed Routing}
\label{app:trace-dotti}

\newtcolorbox{llmcall}[2][]{%
  enhanced jigsaw, breakable,
  colback=blue!3, colframe=blue!55!black,
  fonttitle=\bfseries\small, coltitle=white,
  title={#2}, left=4pt, right=4pt, top=3pt, bottom=3pt,
  boxrule=0.5pt, arc=2pt,
  breakable=true, pad at break*=2mm,
  before upper={\parindent=0pt}, #1}
\newtcolorbox{opstate}[1][]{%
  enhanced jigsaw, breakable,
  colback=gray!8, colframe=gray!50,
  left=4pt, right=4pt, top=2pt, bottom=2pt,
  boxrule=0.3pt, arc=1pt, #1}

\noindent We trace $Q{=}$\emph{``What did the spouse of Andrea Dotti win a
Tony for?''} (gold: \emph{for Best Performance by a Leading Actress in a
Play in Ondine}; MuSiQue~\cite{Trivedi2022}, QID
\texttt{2hop\_\_280451\_84616}) through every operator of
Algorithm~\ref{alg:thirag}. Five LLM-call boxes (in blue) carry the full
prompts and outputs; gray boxes carry the deterministic operator state
(entity-linking, schema projection, coverage, grounding, retrieval) between calls.
At the top level, $\textsc{Covers}$ fails because no
\texttt{dbo:tonyAward} (or any award-category) predicate is instantiated
on the linked entity, so the system refuses to fabricate it and falls
through to symbolic decomposition; $p_1$ is then resolved by SPARQL
(\texttt{dbo:spouse} is covered) and $p_2$ by text retrieval.

\begin{opstate}
\textbf{Steps 1--3 (deterministic, no LLM).}\\
$E{=}\{\texttt{dbr:Andrea\_Dotti\_(psychiatrist)}\}$ (\texttt{dbo:Person});
$\mathcal{S}{=}\{(\texttt{dbo:spouse},\texttt{dbo:Person}),(\texttt{dbo:birthDate},\textit{xsd:date}),\dots\}$;
$\textsc{Covers}(\mathcal{S},Q){=}\textbf{false}$ (\emph{Tony}/\emph{award} unmatched, no \texttt{dbo:tonyAward})
$\Rightarrow$ decompose.
\end{opstate}

\begin{llmcall}{LLM Call \#1 --- Decompose}
\textit{Prompt (abbrev.):}
\begin{small}\begin{verbatim}
Decompose this multi-hop question into independent sub-questions.
QUESTION: What did the spouse of Andrea Dotti win a Tony for?
ENTITIES: Andrea Dotti
SCHEMA:   dbo:Person -> spouse, birthDate, birthPlace, ...
          Actual: dbo:spouse -> dbr:Audrey_Hepburn
Rules:    one fact per sub-Q; mark dependencies;
          bridge questions allowed.
JSON only, no explanation.
\end{verbatim}
\end{small}
\textit{Output:} \verb|{"sub_questions":["Who is the spouse of Andrea Dotti?",|
\verb|"What did [Q0] win a Tony for?"], "dependencies":[[1,0]]}|
\end{llmcall}

\begin{opstate}
\textbf{Resolve $p_1$ ``Who is the spouse of Andrea Dotti?'' --- SPARQL path.}
$\hat{p}_1{=}p_1$ (no predecessor);
$E_{p_1}{=}\{\texttt{dbr:Andrea\_Dotti\_(psychiatrist)}\}$;
$\textsc{Covers}(\mathcal{S}_{p_1},\hat{p}_1){=}\textbf{true}$ (\texttt{dbo:spouse} in Person schema).
\end{opstate}

\begin{llmcall}{LLM Call \#2 --- SPARQL Generation}
\textit{Prompt (abbrev.):}
\begin{small}\begin{verbatim}
You are a SPARQL expert for DBpedia.
Entity:   Andrea Dotti
URI:      <http://dbpedia.org/resource/Andrea_Dotti_(psychiatrist)>
Class:    dbo:Person
  dbo:    spouse, birthDate, birthPlace, ...
  Actual: dbo:spouse -> dbr:Audrey_Hepburn
QUESTION: Who is the spouse of Andrea Dotti?
RULES:    use exact URIs; always add LIMIT 100.
\end{verbatim}
\end{small}
\textit{Output:}
\begin{small}\begin{verbatim}
SELECT ?s WHERE {
  <.../Andrea_Dotti_(psychiatrist)> dbo:spouse ?s . } LIMIT 100
\end{verbatim}
\end{small}
\textit{Execution result:} \verb|[{"s":"dbr:Audrey_Hepburn"}]|
\end{llmcall}

\begin{llmcall}{LLM Call \#3 --- Answer Extraction}
\textit{Prompt:} \verb|Question: Who is the spouse of Andrea Dotti?|\\
\verb|SPARQL Results: [{"s":"Audrey Hepburn"}]|\\
\verb|End your response with FINAL ANSWER: <your answer>|\\
\textit{Output:} \verb|FINAL ANSWER: Audrey Hepburn|
$\Rightarrow M[p_1]\gets\emph{Audrey Hepburn}$.
\end{llmcall}

\begin{opstate}
\textbf{Resolve $p_2$ --- Ground + re-project.}
$\hat{p}_2{=}$``What did Audrey Hepburn win a Tony for?'';
$E_{p_2}{=}\{\texttt{dbr:Audrey\_Hepburn}\}$ (Person);
$\mathcal{S}_{p_2}$: no \texttt{dbo:tonyAward}/award-category property;
$\textsc{Covers}{=}\textbf{false}$
$\Rightarrow$ \textbf{route to text}.
RAG: bi-encoder top-20 $\to$ cross-encoder top-8.
Reranked chunk: \emph{``Audrey Hepburn\dots won the Tony Award for
Best Performance by a Leading Actress in a Play for her role in
Ondine\dots''}
\end{opstate}

\begin{llmcall}{LLM Call \#4 --- Text Reasoning}
\textit{Prompt (abbrev.):}
\begin{small}\begin{verbatim}
Sub-question: What did Audrey Hepburn win a Tony for?
Known facts: Q0 = Audrey Hepburn
Wikipedia: === Audrey Hepburn / (Lead) === ... won the Tony Award
for Best Performance by a Leading Actress in a Play for her role
in Ondine ...
Instructions: brief factual answer; end with FINAL ANSWER: <answer>
\end{verbatim}
\end{small}
\textit{Output:} \verb|FINAL ANSWER: Best Performance by a|
\verb|Leading Actress in a Play in Ondine|
$\Rightarrow M[p_2]\gets\emph{Best Performance by a Leading Actress in a Play in Ondine}$.
\end{llmcall}

\begin{llmcall}{LLM Call \#5 --- Synthesize}
\textit{Prompt (abbrev.):}
\begin{small}\begin{verbatim}
Give ONLY the direct answer, most concise way possible.
Question: What did the spouse of Andrea Dotti win a Tony for?
Information Summary:
- [SPARQL] Who is the spouse of Andrea Dotti? -> Audrey Hepburn
- [TEXT]   What did [Q0] win a Tony for? ->
           Best Performance by a Leading Actress in a Play in Ondine
FINAL ANSWER: <your answer>
\end{verbatim}
\end{small}
\textit{Output:} \verb|FINAL ANSWER: for Best Performance by a Leading|
\verb|Actress in a Play in Ondine| $\checkmark$
\end{llmcall}

\begin{table}[h]
\centering
\caption{I/O summary for the Andrea Dotti trace (cross-reference for the boxes above).}
\label{tab:trace-dotti-io-summary-full}
\footnotesize
\setlength{\tabcolsep}{3pt}
\resizebox{\linewidth}{!}{%
\begin{tabular}{@{}c l l l c@{}}
\toprule
\textbf{\#} & \textbf{Operator} & \textbf{Input} & \textbf{Output} & \textbf{LLM} \\
\midrule
1  & EntityLink     & $Q$                  & \texttt{dbr:Andrea\_Dotti\_(psychiatrist)}                  & ---  \\
2  & SchemaProject  & entity + $Q$         & \texttt{spouse}$\checkmark$, \texttt{tonyAward}$\times$    & ---  \\
3  & Covers         & $\mathcal{S}$        & \textbf{false}                                              & ---  \\
4  & Decompose      & $Q$, schema          & $p_1, p_2$; $p_2{\to}p_1$                                   & \#1  \\
\midrule
5  & Ground($p_1$)  & $M{=}\emptyset$      & ``Who is the spouse of Andrea Dotti?''                      & ---  \\
6  & Covers($p_1$)  & Person schema        & \textbf{true}                                              & ---  \\
7  & Text2SPARQL    & schema + $\hat{p}_1$ & SPARQL query                                                & \#2  \\
8  & Execute        & query                & \texttt{dbr:Audrey\_Hepburn}                                & ---  \\
9  & ExtractAnswer  & results              & ``Audrey Hepburn''                                          & \#3  \\
\midrule
10 & Ground($p_2$)  & $M[p_1]$             & ``What did Audrey Hepburn win a Tony for?''                 & ---  \\
11 & Covers($p_2$)  & Person schema        & \textbf{false}                                             & ---  \\
12 & RAG retrieve   & $\hat{p}_2$, corpus  & Wikipedia lead of Audrey Hepburn                            & ---  \\
13 & TextReason     & chunks + $\hat{p}_2$ & ``Best Performance\dots Ondine''                            & \#4  \\
\midrule
14 & Synthesize     & $M{=}\{p_1, p_2\}$   & \textbf{for Best Performance\dots Ondine}                   & \#5  \\
\bottomrule
\end{tabular}
}% end \resizebox
\end{table}

\noindent 5 LLM calls total. The routing decision (SPARQL for $p_1$, text for $p_2$)
is determined entirely by symbolic schema coverage---no learned classifier.

% ============================================================================
% Deprecated 2026-04-29: Atomised trace below replaced by Dotti above.
% Kept inside \iffalse per project convention; do not delete.
% See app:trace-dotti for active trace.
% ============================================================================
\iffalse

% ============================================================================
% Active trace: compact tcolorbox-per-LLM-call format (5 calls + state blocks).
% Q = "Who is the mother of the director of film Atomised (Film)?"
% Gold: Gisela Elsner | Predicted: Gisela Elsner (correct)
% From real evaluation run (hybrid_eval_1000.jsonl)
% ============================================================================

\subsection{End-to-End Trace with Full LLM Prompts for Mixed Routing}
\label{app:trace-atomised}

\newtcolorbox{llmcall}[2][]{%
  enhanced jigsaw, breakable,
  colback=blue!3, colframe=blue!55!black,
  fonttitle=\bfseries\small, coltitle=white,
  title={#2}, left=4pt, right=4pt, top=3pt, bottom=3pt,
  boxrule=0.5pt, arc=2pt,
  breakable=true, pad at break*=2mm,
  before upper={\parindent=0pt}, #1}
\newtcolorbox{opstate}[1][]{%
  enhanced jigsaw, breakable,
  colback=gray!8, colframe=gray!50,
  left=4pt, right=4pt, top=2pt, bottom=2pt,
  boxrule=0.3pt, arc=1pt, #1}

\noindent We trace $Q{=}$\emph{``Who is the mother of the director of film
Atomised (Film)?''} (gold: \emph{Gisela Elsner}) through every operator of
Algorithm~\ref{alg:thirag}. Five LLM-call boxes (in blue) carry the full
prompts and outputs; gray boxes carry the deterministic operator state
(entity-linking, schema projection, coverage, grounding, retrieval) between calls.

\begin{opstate}
\textbf{Steps 1--3 (deterministic, no LLM).}\\
$E = \{\texttt{dbr:Atomised\_(film)}\}$;
$\mathcal{S}$: \texttt{director}$\checkmark$, \texttt{mother}/\texttt{parent}$\times$;
$\textsc{Covers}(\mathcal{S},Q){=}\textbf{false}\Rightarrow$ decompose.
\end{opstate}

\begin{llmcall}{LLM Call \#1 --- Decompose}
\textit{Prompt (abbrev.):}
\begin{small}\begin{verbatim}
Decompose this multi-hop question into independent sub-questions.
QUESTION: Who is the mother of the director of film
          Atomised (Film)?
ENTITIES: Atomised
SCHEMA:   dbo:Film -> director, producer, ...
          Actual: dbo:director -> Oskar Roehler
Rules:    one fact per sub-Q; mark dependencies;
          bridge questions allowed.
JSON only, no explanation.
\end{verbatim}
\end{small}
\textit{Output:} \verb|{"sub_questions":["Who directed Atomised?",|
\verb|"Who is the mother of [Q0]?"], "dependencies":[[1,0]]}|
\end{llmcall}

\begin{opstate}
\textbf{Resolve $p_1$ ``Who directed Atomised?'' --- SPARQL path.}
$\textsc{Covers}(\mathcal{S}_{p_1},\hat{p}_1){=}\textbf{true}$ (\texttt{director} in Film schema).
\end{opstate}

\begin{llmcall}{LLM Call \#2 --- SPARQL Generation}
\textit{Prompt (abbrev.):}
\begin{small}\begin{verbatim}
You are a SPARQL expert for DBpedia.
Entity:   Atomised
URI:      <http://dbpedia.org/resource/Atomised_(film)>
Class:    dbo:Film
  dbo:    director, producer, starring, ...
  Actual: dbo:director -> Oskar Roehler
QUESTION: Who directed Atomised?
RULES:    use exact URIs; always add LIMIT 100.
\end{verbatim}
\end{small}
\textit{Output:}
\begin{small}\begin{verbatim}
SELECT ?director WHERE {
  <.../Atomised_(film)> dbo:director ?director . } LIMIT 100
\end{verbatim}
\end{small}
\textit{Execution result:} \verb|[{"director":"dbr:Oskar_Roehler"}]|
\end{llmcall}

\begin{llmcall}{LLM Call \#3 --- Answer Extraction}
\textit{Prompt:} \verb|Question: Who directed Atomised?|\\
\verb|SPARQL Results: [{"director":"Oskar Roehler"}]|\\
\verb|End your response with FINAL ANSWER: <your answer>|\\
\textit{Output:} \verb|FINAL ANSWER: Oskar Roehler|
$\Rightarrow M[p_1]\gets\emph{Oskar Roehler}$.
\end{llmcall}

\begin{opstate}
\textbf{Resolve $p_2$ --- Ground + re-project.}
$\hat{p}_2{=}$``Who is the mother of Oskar Roehler?'';
$E_{p_2}{=}\{\texttt{dbr:Oskar\_Roehler}\}$ (Person);
$\textsc{Covers}{=}\textbf{false}$ (no \texttt{mother}/\texttt{parent})
$\Rightarrow$ \textbf{route to text}.
RAG: bi-encoder top-20 $\to$ cross-encoder top-8.
Reranked chunk: \emph{``Oskar Roehler\dots was born in Starnberg as the son of
writer Gisela Elsner and the writer Klaus Roehler.''}
\end{opstate}

\begin{llmcall}{LLM Call \#4 --- Text Reasoning}
\textit{Prompt (abbrev.):}
\begin{small}\begin{verbatim}
Sub-question: Who is the mother of Oskar Roehler?
Known facts: Q0 = Oskar Roehler
Wikipedia: === Oskar Roehler / (Lead) === ... son of writer Gisela
Elsner and the writer Klaus Roehler. ...
Instructions: brief factual answer; end with FINAL ANSWER: <answer>
\end{verbatim}
\end{small}
\textit{Output:} \verb|FINAL ANSWER: Gisela Elsner|
$\Rightarrow M[p_2]\gets\emph{Gisela Elsner}$.
\end{llmcall}

\begin{llmcall}{LLM Call \#5 --- Synthesize}
\textit{Prompt (abbrev.):}
\begin{small}\begin{verbatim}
Give ONLY the direct answer, most concise way possible.
Question: Who is the mother of the director of film Atomised (Film)?
Information Summary:
- [SPARQL] Who directed Atomised? -> Oskar Roehler
- [TEXT]   Who is the mother of [Q0]? -> Gisela Elsner
FINAL ANSWER: <your answer>
\end{verbatim}
\end{small}
\textit{Output:} \verb|FINAL ANSWER: Gisela Elsner| $\checkmark$
\end{llmcall}

\begin{table}[h]
\centering
\caption{I/O summary for the Atomised trace (cross-reference for the boxes above).}
\label{tab:trace-io-summary-full}
\footnotesize
\setlength{\tabcolsep}{3pt}
\begin{tabular}{@{}c l l l c@{}}
\toprule
\textbf{\#} & \textbf{Operator} & \textbf{Input} & \textbf{Output} & \textbf{LLM} \\
\midrule
1  & EntityLink     & $Q$                  & \texttt{dbr:Atomised\_(film)}                       & ---  \\
2  & SchemaProject  & entity + $Q$         & \texttt{director}$\checkmark$, \texttt{mother}$\times$ & ---  \\
3  & Covers         & $\mathcal{S}$        & \textbf{false}                                       & ---  \\
4  & Decompose      & $Q$, schema          & $p_1, p_2$; $p_2{\to}p_1$                            & \#1  \\
\midrule
5  & Ground($p_1$)  & $M{=}\emptyset$      & ``Who directed Atomised?''                            & ---  \\
6  & Covers($p_1$)  & Film schema          & \textbf{true}                                        & ---  \\
7  & Text2SPARQL    & schema + $\hat{p}_1$ & SPARQL query                                          & \#2  \\
8  & Execute        & query                & \texttt{dbr:Oskar\_Roehler}                          & ---  \\
9  & ExtractAnswer  & results              & ``Oskar Roehler''                                    & \#3  \\
\midrule
10 & Ground($p_2$)  & $M[p_1]$             & ``mother of Oskar Roehler?''                          & ---  \\
11 & Covers($p_2$)  & Person schema        & \textbf{false}                                       & ---  \\
12 & RAG retrieve   & $\hat{p}_2$, corpus  & Wikipedia lead of Oskar Roehler                       & ---  \\
13 & TextReason     & chunks + $\hat{p}_2$ & ``Gisela Elsner''                                    & \#4  \\
\midrule
14 & Synthesize     & $M{=}\{p_1, p_2\}$   & \textbf{Gisela Elsner}                               & \#5  \\
\bottomrule
\end{tabular}
\end{table}

\noindent 5 LLM calls total. The routing decision (SPARQL for $p_1$, text for $p_2$)
is determined entirely by schema coverage---no learned classifier.

\fi % --- end deprecated Atomised block (replaced by app:trace-dotti) ---

% =============================================================================
% Error Analysis (moved from main paper input/error-analysis.tex)
% =============================================================================
% Trimmed from the LNCS submission (no main-text \ref to sec:error-analysis,
% tab:error-clusters, or tab:hop-depth). Preserved here in the dbpedia_qa
% codebase as supplementary material. Cross-references (Algorithm~\ref{alg:thirag},
% Table~\ref{tab:trace}, citations) point into the main paper PDF.
% =============================================================================

% FACTCHECK STATUS: PASS (2026-04-29) — all numerical claims verified against
% validated failure-analysis report (project_failure_analysis_run1.md), which
% itself was factchecked PASS against raw judged JSONL traces (judge: gemma3:12b).
% No discrepancies found. Safe to submit.
\section{Error Analysis}
\label{sec:error-analysis}

To understand \emph{why} \sysname\ fails, we classify the $589$
judged-incorrect questions across the three benchmarks
($106$/$1000$ on 2Wiki, $139$/$1000$ on HotpotQA, $344$/$1000$ on
MuSiQue; aggregate fail rate $19.6\%$) by inspecting their full
pipeline traces (routing decisions, sub-question logs, retrieved
chunks). All counts below are verified against the raw judged JSONL
(judge: \textsc{gemma3:12b}). MuSiQue's failure rate ($34.4\%$) is
$3.2\times$ the 2Wiki rate, consistent with its higher hop count
and lower KG connectivity (Table~\ref{tab:hop-depth}); we do not
attempt to fully decouple hop depth from other inter-benchmark
construction differences.

\smallskip\noindent\textbf{Failure taxonomy.} Table~\ref{tab:error-clusters}
groups the $589$ failures into ten clusters. Six dominate ($\geq 5\%$
of failures); we describe each below.

\begin{table}[t]
\centering
\caption{Failure clusters across $589$ judged-incorrect questions;
counts per benchmark and share of total.}
\label{tab:error-clusters}
\small
\begin{tabular}{clrrrrr}
\toprule
\# & Cluster & 2Wiki & HPQA & MuSiQue & Total & \% \\
\midrule
I    & KG forces text, text fails       & 61 & 74 & 102 & 237 & 40.2 \\
II   & Sub-question returns wrong fact  & 2  & 6  & 82  & 90  & 15.3 \\
III  & Routing bypass on multi-hop      & 22 & 30 & 13  & 65  & 11.0 \\
IV   & Under-decomposition              & 3  & 0  & 85  & 88  & 14.9 \\
V    & Over-decomposition               & --- & --- & 38 & 38  & 6.5  \\
VI   & Comparison-logic error           & 18 & 16 & --- & 34  & 5.8  \\
VII  & Entity-link / null sub-Q         & 0  & 8  & --- & 8   & 1.4  \\
VIII & Hedging / refusal                & 0  & 0  & 14  & 14  & 2.4  \\
IX   & DBpedia--Wikipedia format gap    & 0  & 5  & --- & 5   & 0.8  \\
X    & Chain-text-miss intermediate     & 0  & 0  & 10  & 10  & 1.7  \\
\midrule
     & \textbf{Total}                   & 106 & 139 & 344 & \textbf{589} & \textbf{100} \\
\bottomrule
\end{tabular}
\end{table}

\smallskip\noindent\textbf{Cluster I: KG forces text, text fails
($40.2\%$, the dominant mode).}
The KG signals absence of a 2-hop path; the question falls back to
text retrieval, which then fails for benchmark-specific reasons. The
trigger is shared, but the downstream failure differs:
\emph{(i)} on 2Wiki ($61$ cases), the system retrieves the right
Wikipedia page but the predicted string refers to the same concept
under a different label (e.g.\ ``Ottoman Empire'' vs.\ ``Turkey''),
so EM/F1 string-match fails;
\emph{(ii)} on HotpotQA ($74$ cases), the BM25 + reranker pipeline
promotes a topically similar but wrong Wikipedia page (e.g.\
the director's biography instead of the film page), all $74$
exhibiting topically-related but incorrect retrieval;
\emph{(iii)} on MuSiQue ($102$ cases), DBpedia connectivity collapses
beyond two hops (Table~\ref{tab:hop-depth}), so text must traverse
the entire chain unaided, which Wikipedia pages do not support.

\smallskip\noindent\textbf{Cluster IV: under-decomposition
($14.9\%$, almost entirely MuSiQue).}
The decomposition LLM produces fewer sub-questions than the
hop count: $46$ cases of $4\text{-hop}\!\to\!3$ sub-Qs, $30$ of
$3\text{-hop}\!\to\!2$, $7$ of $4\text{-hop}\!\to\!2$, $2$ of
$2\text{-hop}\!\to\!1$. About $80\%$ of these $88$ cases co-occur with all-text retrieval:
when the KG cannot help, decomposition errors compound. We leave
prompt revisions targeting this pattern to future work.

\smallskip\noindent\textbf{Cluster II: sub-question returns wrong
fact ($15.3\%$, dominant on MuSiQue).}
A sub-question executes but retrieves the wrong answer. We focus
on MuSiQue's $82$ cases, where the pattern is sharpest; the
remaining $8$ ($2$ on 2Wiki, $6$ on HotpotQA) were not separately
re-classified. On the $82$ MuSiQue cases, $91.5\%$ ($75$/$82$) have
the gold entity \emph{absent from the entire retrieved chain}---only
$7$/$82$ are genuine final-synthesis errors, suggesting the
bottleneck is retrieval, not synthesis: improvements to
the synthesis prompt will not help if the gold fact was never
retrieved.

\smallskip\noindent\textbf{Cluster III: routing bypass
($11.0\%$, universal across benchmarks).}
Multi-hop questions are answered by SPARQL alone because routing
matches a surface predicate (e.g.\ ``director'' $\to$
\texttt{dbo:director}) and returns early, skipping decomposition
($22$/$30$/$13$ on the three benchmarks). A single guard---force
decomposition consideration when the question is multi-hop---would,
in principle, unblock decomposition for all $65$ cases; the recoverable subset is bounded
by sub-question success rates and is not directly measured here.

\smallskip\noindent\textbf{Cluster V: over-decomposition
($6.5\%$, MuSiQue only).}
The opposite failure of Cluster IV: $14$/$38$ cases follow a
$2$-hop $\!\to\!3{+}$ sub-question pattern in which a phantom sub-question
introduces retrieval noise that propagates into the final answer;
a further $5$--$7$ cases are $4$-hop$3$ questions split into $5$--$6$
sub-questions, and the remaining cases are over-decomposed by other
patterns we do not separately enumerate. The fix points to
bidirectional decomposition guidance (produce roughly $N$ sub-questions
for an $N$-hop question), not just a stricter floor.

\smallskip\noindent\textbf{Cluster VI: comparison logic
($5.8\%$).}
Apparent reasoning errors that on inspection are date/string-format
or value-multiplicity mismatches---for example, a film with multiple
listed release dates ($1948$, $1968$, $1991$) compared against a
single-date counterpart ($1988$) causes the comparison prompt to
flip the answer; nationality-equivalence cases (e.g.\ ``Greek and
Italian'' vs.\ ``Greek'') exhibit a similar pattern. A canonical
date and nationality mapper applied prior to comparison addresses
these cases.

\smallskip\noindent\textbf{Hop-depth as a structural ceiling.}
Table~\ref{tab:hop-depth} shows the joint behaviour of failure rate
and DBpedia chain connectivity on MuSiQue. Connectivity drops
sharply with hop count and the failure rate climbs accordingly;
at $4$-hop with intermediate bridge type ``$4$hop$2$'', connectivity
reaches $0\%$ and the system fails on $92.6\%$ of questions
($n=27$, small sample). Past two hops, DBpedia connectivity collapses
on these benchmarks; prompt-level changes are unlikely to recover this
regime without KG expansion or stronger external retrieval.

\begin{table}[t]
\centering
\caption{MuSiQue failure rate vs.\ DBpedia chain connectivity by
hop type. Connectivity = fraction of questions whose entity chain
is fully traversable in DBpedia.}
\label{tab:hop-depth}
\small
\begin{tabular}{lrr}
\toprule
Hop type & Fail \% & DBpedia connectivity \\
\midrule
2hop   & 28.6 & 32.8 \\
3hop1  & 24.7 & 14.4 \\
3hop2  & 45.2 &  1.4 \\
4hop1  & 55.6 &  2.8 \\
4hop2  & \textbf{92.6} &  0.0 \\
4hop3  & 58.1 &  0.0 \\
\bottomrule
\end{tabular}
\end{table}

\smallskip\noindent\textbf{What is \emph{not} the bottleneck.}
Two mechanisms are not the bottleneck---entity linking and final-synthesis
reasoning---even though both could plausibly be.
Of the $8$ HotpotQA Cluster~VII failures, none stem from entity
linking; all are null sub-question outputs from downstream LLM
failure. Entity linking is also absent from 2Wiki and MuSiQue top
causes; the Spotlight + SPARQL fallback chain does not appear
among the top failure causes.
Final-synthesis reasoning is also a minor contributor: only $7$/$82$
MuSiQue Cluster II cases are true synthesis errors.

\smallskip\noindent\textbf{Caveats.}
All counts use \textsc{gemma3:12b} as judge, which is more lenient
than strict EM; a stricter judge would shift Cluster II and VI counts
upward. The $4$hop$2$ regime ($n=27$) is small; we treat the $92.6\%$
figure as indicative rather than conclusive. One untracked failure
mode---final-synthesis hallucination overriding a literal
``\texttt{Unknown}'' sub-question answer with parametric memory---is
currently absorbed into Cluster~II ($\geq 1$ confirmed; up to $16$
MuSiQue cases involve sub-questions returning literal
``\texttt{Unknown}'' where override is plausible but unverified).

% OLD (flat list, "headline" reads as boastful):
% \smallskip\noindent\textit{In summary,} four findings drive future
% work: (1) $40\%$ of failures arise when the KG forces a text
% fallback that fails for label-, retrieval-, or chain-disconnect
% reasons---the headline failure mode; (2) decomposition errors are
% bidirectional, with both under- and over-decomposition observable;
% (3) routing bypass on multi-hop questions is the cheapest universal
% fix ($65$ cases, single guard); (4) on MuSiQue, hop depth dictates
% outcome through DBpedia connectivity, defining a structural ceiling
% that prompt-level changes alone cannot lift without KG expansion.
\smallskip\noindent\textit{In summary,} four findings drive
future work, paired by what they implicate. The KG draws a
boundary: (1) $40\%$ of failures arise when the KG forces a
text fallback that fails for label-, retrieval-, or
chain-disconnect reasons; (4) on
MuSiQue, hop depth dictates outcome through DBpedia
connectivity, defining a ceiling that prompt-level changes
alone cannot lift without KG expansion. Routing leaks inside
that boundary: (2) decomposition errors are bidirectional,
with both under- and over-decomposition observable; (3)
routing bypass on multi-hop questions is the cheapest
universal fix ($65$ cases, single guard).
